% ===== PRÉAMBULE CANONIQUE — à coller VERBATIM en tête de CHAQUE .tex =====
\documentclass[11pt,a4paper]{article}
\usepackage[utf8]{inputenc}\usepackage[T1]{fontenc}\usepackage{lmodern}
\usepackage[french]{babel}
\usepackage{amsmath,amssymb,mathtools}\usepackage{icomma}
\usepackage{siunitx}\sisetup{locale=FR}  % \qty{}{}, \si{}, \ang{}
\usepackage{xcolor}\usepackage{colortbl}
\usepackage{tikz}\usepackage{pgfplots}\pgfplotsset{compat=1.18}
\usepackage[siunitx,europeanresistors]{circuitikz} % circuits (élec.)
\usepackage[most]{tcolorbox}\usepackage{enumitem}\usepackage{array,booktabs}
\usepackage[a4paper,margin=2.2cm]{geometry}
\usetikzlibrary{arrows.meta,calc,angles,quotes,positioning,patterns,%
  backgrounds,shapes.geometric,decorations.pathmorphing,decorations.markings,intersections}
\definecolor{primary}{RGB}{222,49,99}\definecolor{primarydark}{RGB}{150,22,55}
\definecolor{defcol}{RGB}{0,114,178}\definecolor{methcol}{RGB}{52,92,148}
\definecolor{excol}{RGB}{0,135,90}\definecolor{attncol}{RGB}{200,40,55}
\tcbset{boxbase/.style={enhanced,breakable,boxrule=0.9pt,arc=2.5pt,
  left=3mm,right=3mm,top=2.3mm,bottom=2.3mm,fonttitle=\bfseries,coltitle=white}}
% --- boîtes TUTORIEL (noms EXACTS, ne pas renommer) ---
\newtcolorbox{cle}[1][]{boxbase,colback=defcol!6,colframe=defcol,
  title={\textsc{À retenir}\ifx\relax#1\relax\relax\else\textmd{ : \itshape #1}\fi}}
\newtcolorbox{methode}[1][]{boxbase,colback=methcol!7,colframe=methcol,sharp corners,
  title={\textsc{Méthode}\ifx\relax#1\relax\relax\else\textmd{ --- \itshape #1}\fi}}
\newtcolorbox{exemplebox}[1][]{boxbase,colback=excol!5,colframe=excol,
  title={\textsc{Exemple}\ifx\relax#1\relax\relax\else\textmd{ : \itshape #1}\fi}}
\newtcolorbox{piege}[1][]{boxbase,colback=attncol!6,colframe=attncol,
  title={\textsc{Piège}\ifx\relax#1\relax\relax\else\textmd{ : \itshape #1}\fi}}
\newtcolorbox{remarque}[1][]{boxbase,colback=black!4,colframe=black!45,
  title={\textsc{Remarque}\ifx\relax#1\relax\relax\else\textmd{ : \itshape #1}\fi}}
% --- boîtes EXERCICE / CORRIGÉ (noms EXACTS, auto-comptées) ---
\newtcolorbox[auto counter]{exo}[1][]{boxbase,colback=primary!4,colframe=primary,
  title={\textsc{Exercice}~\thetcbcounter\ifx\relax#1\relax\relax\else\textmd{ --- \itshape #1}\fi}}
\newtcolorbox[auto counter]{corrige}[1][]{boxbase,colback=excol!4,colframe=excol,
  title={\textsc{Corrigé}~\thetcbcounter\ifx\relax#1\relax\relax\else\textmd{ --- \itshape #1}\fi}}
\newcommand{\vv}[1]{\vec{#1}}\newcommand{\ds}{\displaystyle}
\newcommand{\R}{\mathbb{R}}\newcommand{\N}{\mathbb{N}}\newcommand{\Z}{\mathbb{Z}}
% (un \usepackage{fancyhdr}+en-tête est possible ; il est ignoré côté web)
% =========================================================================
\begin{document}

\begin{center}
{\large\bfseries\color{excol} Thème 4 — Corrigé Difficile}\\[-2pt]
{\color{excol}\rule{5cm}{1pt}}
\end{center}

\begin{corrige}[un point sur deux sources]
La différence de marche décide de tout : $\delta=n\lambda$ (constructif) ou
$(2n+1)\lambda/2$ (destructif).
\begin{enumerate}[label=\alph*)]
  \item $\delta=d_1-d_2=10-7=\SI{3}{\centi\metre}.$
  \item $\dfrac{\delta}{\lambda}=\dfrac{3}{2}=1{,}5$ (demi-entier), soit
        $\delta=(2\times1+1)\lambda/2=3\lambda/2$ : interférence \textbf{destructive}.
  \item Pour rendre $M$ constructif, il faut $\delta=n\lambda'$. Avec $n=1$ :
        $\lambda'=\delta=\SI{3}{\centi\metre}$ (ou $\lambda'=\SI{1.5}{\centi\metre}$ pour $n=2$,
        etc.).
\end{enumerate}
\end{corrige}

\begin{corrige}[combien de franges brillantes sur l'écran ?]
\begin{enumerate}[label=\alph*)]
  \item $i=\dfrac{\lambda D}{a}=\dfrac{\num{6e-7}\times 2}{\num{5e-4}}
        =\num{2.4e-3}\,\si{\metre}=\SI{2.4}{\milli\metre}.$
  \item On cherche les entiers $n$ tels que $|n\,i|\le\SI{10}{\milli\metre}$, soit
        $|n|\le\dfrac{10}{2{,}4}\approx 4{,}17$. Donc $n$ va de $-4$ à $+4$ : cela fait
        $9$ valeurs $\Rightarrow$ \textbf{9 franges brillantes} (la centrale $n=0$ comprise).
\end{enumerate}
\end{corrige}

\begin{corrige}[deux haut-parleurs et un silence]
\begin{enumerate}[label=\alph*)]
  \item $\lambda=\dfrac{v}{f}=\dfrac{340}{1700}=\SI{0.2}{\metre}=\SI{20}{\centi\metre}.$
  \item \textbf{Oui.} Un silence se produit là où l'interférence est destructive :
        $\delta=(2n+1)\lambda/2$. La plus petite valeur est
        $\delta=\dfrac{\lambda}{2}=\SI{0.1}{\metre}=\SI{10}{\centi\metre}.$
  \item En ce point, les deux sons arrivent \emph{en opposition de phase} : ils s'annulent
        mutuellement, d'où le silence.
\end{enumerate}
\end{corrige}

\begin{corrige}[deux ondes en quadrature]
\begin{enumerate}[label=\alph*)]
  \item Avance de $\lambda/4$ : $\Delta\varphi=2\pi\,\dfrac{\lambda/4}{\lambda}=\dfrac{\pi}{2}\ \si{\radian}.$
  \item $\Delta\varphi=\pi/2$ n'est ni $0$ (qui donnerait $2X$) ni $\pi$ (qui donnerait $0$) :
        l'amplitude résultante a une \textbf{valeur intermédiaire}, comprise entre $X$ et $2X$.
  \item $A=2X\cos(\Delta\varphi/2)=2X\cos(\pi/4)=2X\times\dfrac{\sqrt2}{2}=X\sqrt2\approx 1{,}41\,X.$
\end{enumerate}
\end{corrige}

\begin{corrige}[rouge ou bleu : quelles franges plus larges ?]
$i=\lambda D/a$ : à géométrie fixée, $i$ est \emph{proportionnel} à $\lambda$.
\begin{enumerate}[label=\alph*)]
  \item Le rouge a la plus grande longueur d'onde ($\lambda_R>\lambda_B$), donc il donne
        l'\textbf{interfrange le plus grand}.
  \item $i_R=\dfrac{\num{7e-7}\times 2}{\num{4e-4}}=\SI{3.5}{\milli\metre}$ ;
        $i_B=\dfrac{\num{4.5e-7}\times 2}{\num{4e-4}}=\SI{2.25}{\milli\metre}.$
  \item $\dfrac{i_R}{i_B}=\dfrac{3{,}5}{2{,}25}\approx 1{,}56=\dfrac{700}{450}=\dfrac{\lambda_R}{\lambda_B}$ :
        le rapport des interfranges est bien le rapport des longueurs d'onde.
\end{enumerate}
\end{corrige}

\end{document}
